\documentclass[10pt,a4paper]{amsart}
\usepackage[margin=19mm]{geometry}
\usepackage{amsmath,amssymb,mathtools}
\usepackage[scr=rsfs,cal=euler]{mathalfa}
\usepackage{xfrac}
\usepackage[unicode,hidelinks,bookmarks=false]{hyperref}
\newcommand{\ie}{i.e.\ }
\newcommand{\cf}{cf.\ }
\newcommand{\resp}{respectively\ }
\newcommand{\Subsection}{\subsection}
\providecommand{\nobreakdash}{\nobreak\hbox{-}}
\setlength{\emergencystretch}{2em}
\multlinegap0pt
\usepackage{bm}
\usepackage{bbm}
\usepackage{dsfont}
\usepackage{xspace}
\usepackage{cancel}
\usepackage{mathtools}
\usepackage{cleveref}
\usepackage{amsthm}
\makeatletter
\@ifundefined{lemma}{\newtheorem{lemma}{Lemma}}{}
\makeatother
\def\R{{\mathbb{R}}}
\def\erre{{\mathbb{R}}}
\title[Quantitative stability of the Rossby--Haurwitz waves of degr]{Quantitative stability of the Rossby--Haurwitz waves of degree two for the Euler equation on $\mathbb{S}^2$}
\author{Matias G. Delgadino, Luca Melzi}
\date{Source: arXiv:2509.16156v1 (CC BY 4.0)}
\begin{document}
\maketitle

\noindent\emph{Source and licence.} Quantitative stability of the Rossby--Haurwitz waves of degree two for the Euler equation on $\mathbb{S}^2$. Version arXiv:2509.16156v1 (CC BY 4.0), distributed under the Creative Commons Attribution 4.0 International licence, as recorded in the arXiv licence metadata for this exact version and shown on the abstract page https://arxiv.org/abs/2509.16156v1. Full rights notice: licenses/arxiv-2509.16156-CC-BY-4.0.txt.\par\smallskip

\noindent\textit{Editorial scope.} Counted unit: the Lemma whose source label is \texttt{lemma:ker\_1dim} in the article (source0.tex lines 350--364, source label \texttt{lemma:ker\_1dim}), delivered together with the article's own complete original proof of it (source0.tex lines 409--465, one \texttt{proof} environment of 57 lines). No mathematics is written, repaired or re-derived: statement and proof are the article's own text line for line. The proof is detached from the statement in the source: the article states the result at lines 350--364 and prints its proof at lines 409--465, and the proof environment's own header names the statement, so the association is the article's own. Required context retained uncounted: lemma:grad\_neq0 (lines 338--343); eq:taylor\_expansion (lines 372--374). Every definition, display and label of those lines is retained unchanged. Omitted from this package and not redistributed: 1--337, 344--349, 365--371, 375--408, 466--1054. Nothing inside a retained excerpt is omitted or altered: every retained body is delivered exactly as the article's lines are written, so no omission receipt of any kind accompanies this package. Citations: the retained excerpts carry no citation command at all, so no bibliography entry of the article is transcribed and no reference is redistributed. Reference adaptation: none was needed and none was made; every reference inside the delivered unit resolves inside it, so no unresolved reference or citation is produced. Printed slips kept literal: no printed word, symbol, hypothesis or number of the counted statement or of its proof was altered. Layout and class: the article's own class is \texttt{amsart} with packages including graphicx, subfig, tabularx, transparent, eso-pic, tikz, caption, xcolor, amsthm, thmtools, float, geometry, inputenc, fontenc; the wrapper is the lane's pinned \texttt{amsart} editorial wrapper at 10pt on A4, which restates the theorem-like environments the retained text needs and adds the article's own macro definitions that the retained text needs (\texttt{\textbackslash R}, \texttt{\textbackslash erre}), with extra wrapper packages (bm, bbm, dsfont, xspace, cancel, mathtools, cleveref). The delivered numbering is the unit's own. Assets: no retained span contains a figure environment, a graphics-inclusion command, a PDF-page inclusion, a table or a TikZ picture; no image file is copied into this package, the package declares no asset dependency and no image file is redistributed with it. Licence: version arXiv:2509.16156v1 of this article is distributed under the Creative Commons Attribution 4.0 International licence, as recorded in the arXiv licence metadata for this exact version and shown on the abstract page https://arxiv.org/abs/2509.16156v1. A full rights notice is delivered at licenses/arxiv-2509.16156-CC-BY-4.0.txt. Characters: every retained excerpt is pure ASCII, so the delivered engine, XeLaTeX with its default Latin Modern text font, reports no missing character.
\begin{lemma} \label{lemma:grad_neq0}
        Let $F:\erre^n \to\erre^n$ be a smooth function and $x_0\in \R^n$. Assume that $\det\nabla F(x_0)\neq0$, then there exists $\delta>0$ and $C=C(F,x_0)>0$ such that for all $x\in B_\delta(x_0)$, we have the stability estimate
        \begin{equation}\label{eq:grad_neq0}
            |x-x_0|\leq C |F(x)-F(x_0)|.
        \end{equation}
\end{lemma}

\begin{equation}
    F(x)-F(x_0)=\sum_{k=1}^{N-1}\frac{1}{k!}\nabla^kF(x_0)(x-x_0,\dots,x-x_0)+\frac{1}{N!}\nabla^NF(\xi)(x-x_0,\dots,x-x_0). \label{eq:taylor_expansion}
\end{equation}

\begin{lemma} \label{lemma:ker_1dim}
    Let $F:\erre^n\to\erre^n$ be a smooth function and $x_0\in \R^n$. Assume that
    \begin{equation*}
        \dim\ker\nabla F(x_0)=1,
    \end{equation*}
    and that
    \begin{equation}
        \nabla^2F(x_0)(e,e)\notin \operatorname{col}\nabla F(x_0),
        \label{eq:1dimker_crucial_assumption}
    \end{equation}
    where $e\in\ker\nabla F(x_0)$ is a unit vector. Then, there exists $\delta>0$ and a constant $C=C(F,x_0)>0$, such that for all $x\in B_\delta(x_0)$ we have the stability estimate
    \begin{equation}\label{eq:ker_1dim}
        |x-x_0|^2\leq C|F(x)-F(x_0)|.
    \end{equation}
\end{lemma}

\begin{proof}[Proof of \cref{lemma:ker_1dim}]
    We define
    \begin{equation}
        \lambda:=\inf_{w\in\ker\nabla F(x_0)^\perp}|\nabla F(x_0)w|>0,
    \end{equation}
    and we pick 
    $$
        \gamma=\min\left\{\frac{\lambda}{4\sup_{\xi\in B_1(x_0)}\|\nabla^2 F(\xi)\|},1\right\}.
    $$
    We take $\delta\in(0,\gamma)$ to be chosen later. For any $x\in B_\delta(x_0)$, we decompose $x=x_0+te+sw$ for some $s,t\in[0,\delta)$, with $e\in\ker\nabla F(x_0),\,|e|=1$ and $w\in\ker\nabla F(x_0)^\perp,\,|w|=1$.
    We consider the two cases $\frac{t^2}{\gamma}\leq s$ and $s<\frac{t^2}{\gamma}$.
    
    \textbf{First case.} If $\frac{t^2}{\gamma}\leq s$, the proof goes similarly to the one of \cref{lemma:grad_neq0}.
    Using a Taylor expansion \eqref{eq:taylor_expansion} of order two, taking norms and using the fact that $\nabla F(x_0)e=0$, we have
    \begin{equation*}
        \lambda s-\frac{1}{2}\sup_{\xi\in B_1(x_0)}\|\nabla^2 F(\xi)\| |x-x_0|^2\le\left|\nabla F(x_0)w s+\frac{1}{2}\nabla^2 F(\xi)(x-x_0,x-x_0)\right| = |F(x)-F(x_0)|.
    \end{equation*}
    Using the triangle inequality, and the assumption $\frac{t^2}{\gamma}\leq s$, we have
    \begin{equation*}
        |x-x_0|^2=|sw+te|^2\le 2s^2+2t^2\leq 2(\delta+\gamma) s\leq 4\gamma s,
    \end{equation*}
    where in the last inequality we have used that $\delta<\gamma$. Combining the two previous estimates, we obtain
    \begin{equation*}
        \frac{\lambda}{8\gamma} |x-x_0|^2\le\frac{\lambda}{2} s\le \left(\lambda-2\delta \sup_{\xi\in B_1(x_0)}\|\nabla^2 F(\xi)\|\right) s \le |F(x)-F(x_0)|,
    \end{equation*}
    where we have used the definition of $\gamma$. This shows the desired estimate holds with $C=\frac{8\gamma}{\lambda}$.
    
    \textbf{Second case.} We now consider $s<\frac{t^2}{\gamma}$. We start with the Taylor expansion \eqref{eq:taylor_expansion} of order three, and we take norms to get
    \begin{eqnarray*}
        \underbrace{\left|\nabla F(x_0) w s+\frac{1}{2}\nabla^2 F(x_0)(e,e)t^2+\nabla^2 F(x_0)(e,w)st+\nabla^2 F(x_0)(w,w)t^2+\frac{1}{6}\nabla^3F(\xi)(x-x_0)^3\right|}_{I}\le |F(x)-F(x_0)|.
    \end{eqnarray*}
    Using the triangle inequality, we can estimate
    $$
        \left|\nabla F(x_0) w s+\frac{1}{2}\nabla^2 F(x_0)(e,e)t^2\right|-\sup_{\xi\in B_1(x_0)}\|\nabla^2 F(x_0)\|\left(st+\frac{s^2}{2}\right)-\frac{1}{6}\sup_{\xi\in B_1(x_0)}\|\nabla^3F(\xi)\|(t^3+s^3+3s^2t+3st^2)\le I.
    $$
    We use the assumption of non-degeneracy \eqref{eq:1dimker_crucial_assumption}, and define
    \begin{equation*}
        0<\eta:=|\mathbb{P}_{\operatorname{col}\nabla F(x_0)^\perp}\nabla^2F(x_0)(e,e)|,
    \end{equation*}
    to substitute the first term in the previous inequality, obtaining
    \begin{align*}
        \frac{1}{2}\eta t^2-\sup_{\xi\in B_1(x_0)}\|\nabla^2 F(x_0)\|\left(st+\frac{s^2}{2}\right)-\frac{1}{6}\sup_{\xi\in B_1(x_0)}\nabla^3F(\xi)(t^3+s^3+3s^2t+3st^2)\leq |F(x)-F(x_0)|.
    \end{align*}
    We now use the fact that $s<\frac{t^2}{\gamma}$ and $s,t<\delta<\gamma$, to obtain
    $$
        \left(\frac{1}{2}\eta -\sup_{\xi\in B_1(x_0)}\|\nabla^2 F(x_0)\|\frac{3}{2}\frac{\delta}{\gamma}-\frac{4}{3}\sup_{\xi\in B_1(x_0)}\|\nabla^3F(\xi)\|\delta \right)t^2\leq |F(x)-F(x_0)|.
    $$
    We restrict
    $$
        \delta<\frac{\eta}{4\left(\sup_{\xi\in B_1(x_0)}\|\nabla^2 F(x_0)\|\frac{3}{2\gamma}+\frac{4}{3}\sup_{\xi\in B_1(x_0)}\|\nabla^3F(\xi)\|\right)},
    $$
    to obtain, for all $x\in B_\delta(x_0)$, the estimate
    \begin{equation*}
        \left(3\frac{\delta}{\gamma}+1\right)^{-1}|x-x_0|^2 \le t^2\le \frac{4}{\eta}|F(x)-F(x_0)|.
    \end{equation*}
    Restricting further $\delta<\gamma/3$, we can conclude that the desired estimate \eqref{eq:ker_1dim} holds with $C=\frac{8}{\eta}$.
\end{proof}

\end{document}
